\section*{الفصل \ref{ch:b2:hermitian} --- الصيغ الهرميتية}

\begin{solution}{exo:b2:hermitian:1}
$\langle x, y\rangle = \conj{1}\cdot\iu + \conj{\iu}\cdot 1 = \iu -
\iu = 0$: متعامدتان. $\norm x = \norm y = \sqrt{1 + 1} = \sqrt2$.
والأساس المتعامد الممنظم: $\bigl(\frac{1}{\sqrt2}(1, \iu),\;
\frac{1}{\sqrt2}(\iu, 1)\bigr)$ --- أي الزوج المنظَّم نفسه.
\end{solution}

\begin{solution}{exo:b2:hermitian:2}
الأولى: تساوي منقولها \omterm{def:b2:hermitian:adjoint}{المرافق} (فالقطر حقيقي، والعنصران $\conj{\iu} =
-\iu$ متبادلان): فهي هرميتية (ومن ثم ناظمية)؛ وليست وحدوية (لأن $A^\dagger A
\neq I$: فالأعمدة ليست واحدية).

الثانية: $A^\dagger A = \frac12\begin{pmatrix} 1 & -\iu\\ -\iu &
1\end{pmatrix}\begin{pmatrix} 1 & \iu\\ \iu & 1\end{pmatrix} =
\frac12\begin{pmatrix} 2 & 0\\ 0 & 2\end{pmatrix} = I$: فهي وحدوية
(ومن ثم ناظمية)؛ وليست هرميتية.

الثالثة: $A^\dagger A = E_{22} \neq E_{11} = AA^\dagger$: فهي ليست ناظمية
(ومن ثم ليست هرميتية ولا وحدوية) --- وهي المثال المضاد
القياسي العديم القوى.
\end{solution}

\begin{solution}{exo:b2:hermitian:3}
الوحدانية: يفرض $A = H + \iu K$ مع $H^\dagger = H$، $K^\dagger = K$
أن $A^\dagger = H - \iu K$، ومنه $H = \frac{A + A^\dagger}{2}$،
$K = \frac{A - A^\dagger}{2\iu}$؛ وهاتان الصيغتان هرميتيتان
(للتحقق: $\bigl(\frac{A - A^\dagger}{2\iu}\bigr)^\dagger = \frac{
A^\dagger - A}{-2\iu} = K$) وتعيدان بناء $A$: أي الوجود.

الناظمية: $A^\dagger A - AA^\dagger = (H - \iu K)(H + \iu K) - (H
+ \iu K)(H - \iu K) = 2\iu(HK - KH)$: وهو ينعدم إذا وفقط إذا $HK = KH$.
\end{solution}

\begin{solution}{exo:b2:hermitian:4}
$\chi_A = (X-2)^2 - 1$: القيمتان الذاتيتان $3$ و$1$. والمتجهتان الذاتيتان، بالحساب
المباشر:
\[
A\begin{pmatrix}1\\ -\iu\end{pmatrix}
= \begin{pmatrix} 2 + \iu(-\iu)\\ -\iu + 2(-\iu)\end{pmatrix}
= \begin{pmatrix} 3\\ -3\iu \end{pmatrix}
= 3\begin{pmatrix}1\\ -\iu\end{pmatrix},
\qquad
A\begin{pmatrix}1\\ \iu\end{pmatrix}
= \begin{pmatrix} 2 + \iu\cdot\iu\\ -\iu + 2\iu\end{pmatrix}
= \begin{pmatrix}1\\ \iu\end{pmatrix}.
\]
والأساس الذاتي المتعامد الممنظم: $u_1 = \frac{1}{\sqrt2}(1, -\iu)$
(\omterm{def:b2:reduction:eigen}{بالقيمة الذاتية} $3$)، و$u_2 = \frac{1}{\sqrt2}(1, \iu)$ (\omterm{def:b2:reduction:eigen}{بالقيمة الذاتية}
$1$)؛ والتعامد كما في \cref{exo:b2:hermitian:1}. والقوى، عبر
الإسقاطات الطيفية $A^k = 3^k u_1u_1^\dagger + 1^k\,
u_2u_2^\dagger$:
\[
A^k = U\begin{pmatrix} 3^k & 0\\ 0 & 1\end{pmatrix}U^\dagger
= \frac{3^k}{2}\begin{pmatrix} 1 & \iu\\ -\iu & 1\end{pmatrix}
+ \frac{1}{2}\begin{pmatrix} 1 & -\iu\\ \iu & 1\end{pmatrix}
= \frac12\begin{pmatrix}
3^k + 1 & (3^k - 1)\iu\\
-(3^k-1)\iu & 3^k + 1
\end{pmatrix}.
\]
(وللتحقق من $k = 1$: نستعيد $A$.)
\end{solution}

\begin{solution}{exo:b2:hermitian:5}
متراصة: مغلقة (فهي الصورة العكسية للمجموعة $I$ بالتطبيق \omterm{def:b2:metric:continuity}{المتصل} $U \mapsto
U^\dagger U$) ومحدودة (فالأعمدة متجهات واحدية: بعناصر
طويلتها $\leq 1$) في $\mathcal{M}_n(\C) \simeq \R^{2n^2}$.

\omterm{def:b2:reduction:eigen}{والقيم الذاتية}: $\operatorname{diag}(\lambda, 1, \dots, 1)$ وحدوية
من أجل أي $\abs\lambda = 1$. \omterm{def:b2:linalg:det}{والمحدد}: $\abs{\det U}^2 = \det
U^\dagger \det U = \det(U^\dagger U) = 1$ (باستعمال $\det A^\dagger =
\conj{\det A}$): ومنه $\det U$ يقع على دائرة الوحدة.
\end{solution}

\begin{solution}{exo:b2:hermitian:6}
بالمبرهنة الطيفية، نعمل في أساس ذاتي متعامد ممنظم للمصفوفة $H$:
فيؤول كل شيء إلى أعداد $\lambda \in \R$ (وهي \omterm{def:b2:reduction:eigen}{القيم الذاتية}).
وللمصفوفة $H + \iu I$ \omterm{def:b2:reduction:eigen}{القيم الذاتية} $\lambda + \iu \neq 0$: فهي قابلة للقلب. وللمصفوفة $U$
\omterm{def:b2:reduction:eigen}{القيم الذاتية} $\mu = \frac{\lambda - \iu}{\lambda + \iu}$، وطويلتها
$1$ (لأن $\abs{\lambda - \iu} = \abs{\lambda + \iu}$ من أجل $\lambda$ الحقيقية): ومن ثم يصحّ $U^\dagger U = I$ لأن $U$ قابلة للتقطير
\omterm{def:b2:hermitian:adjoint}{وحدويًا} بقيم ذاتية طويلتها واحد (فهي قطرية في الأساس المتعامد
الممنظم المختار). ولو كان $\mu = 1$ لفُرض $-\iu = \iu$:
وهذا مستحيل، ومنه $1 \notin \operatorname{Sp} U$. (فتحويل كايلي
يرسل الهرميتية إلى الوحدوية ناقص نقطة --- وهي النسخة
المصفوفية للتطبيق من $\R$ إلى الدائرة.)
\end{solution}

\begin{solution}{exo:b2:hermitian:7}
من أجل زوج ذاتي $Ax = \lambda x$ (مع $x \neq 0$): $\lambda\norm x^2 =
\langle x, Ax\rangle > 0$، ومنه $\lambda > 0$ (وهي حقيقية أصلًا،
لأن \cref{prop:b2:hermitian:eigenvalues}). والجذر التربيعي: في أساس
طيفي، $B = U\operatorname{diag}(\sqrt{\lambda_i})U^\dagger$:
فهي هرميتية معرَّفة موجبة، و$B^2 = A$. \omterm{def:b2:linalg:det}{والمحدد}: جداء
\omterm{def:b2:reduction:eigen}{القيم الذاتية} الموجبة.
\end{solution}

\begin{solution}{exo:b2:hermitian:8}
\begin{enumerate}
  \item $\norm{u(x)}^2 = \langle u(x), u(x)\rangle = \langle x,
        u^*u(x)\rangle = \langle x, uu^*(x)\rangle =
        \norm{u^*(x)}^2$. وبتطبيق هذا على التشاكل الناظمي $u -
        \lambda\,\mathrm{id}$ (فمرافقه $u^* -
        \conj\lambda$، والناظمية موروثة): $\norm{(u -
        \lambda)x} = \norm{(u^* - \conj\lambda)x}$، ومن ثم تتطابق
        النواتان.
  \item من أجل متجهتين ذاتيتين $u(x) = \lambda x$، $u(y) = \mu y$
        (مع $\lambda \neq \mu$): باستعمال (1)، $u^*(x) = \conj\lambda x$؛
        ثم
        \[
        \lambda\langle y, x\rangle = \langle y, u(x)\rangle
        = \langle u^*(y), x\rangle = \langle \conj\mu\, y, x\rangle
        = \mu \langle y, x\rangle ,
        \]
        ومنه $\langle y, x\rangle = 0$. واستقرار $E_\lambda^\perp$:
        من أجل $x \perp E_\lambda$ و$z \in E_\lambda$، $\langle z,
        u(x)\rangle = \langle u^*(z), x\rangle =
        \langle\conj\lambda z, x\rangle = 0$.
  \item بالتراجع على البعد: على $\C$، للتشاكل $u$ \omterm{def:b2:reduction:eigen}{متجهة ذاتية}
        $e_1$ (فنُنظّمها)؛ ومتممتها المتعامدة مستقرة
        تحت $u$ (حسب (2)) \emph{وتحت} $u^*$ (بالحجة نفسها
        مع تبادل الدورين)، ومن ثم فالمقصور ناظمي: فنستعمل التراجع
        ونتلاصق الأسس الذاتية المتعامدة الممنظمة. وعكسيًا، فإن التشاكل
        $u = UDU^\dagger$ \omterm{def:b2:reduction:diag}{القابل للتقطير} \omterm{def:b2:hermitian:adjoint}{وحدويًا} يحقق $u^*u
        = U\conj D D U^\dagger = U D\conj D U^\dagger = uu^*$:
        أي إنه ناظمي.
\end{enumerate}
\end{solution}

\begin{solution}{exo:b2:hermitian:9}
\omterm{def:b2:reduction:eigen}{القيم الذاتية}: من أجل $u(x) = \lambda x$، $x \neq 0$:
\[
\lambda\norm x^2 = \langle x, u(x)\rangle
= \langle u^*(x), x\rangle = -\langle u(x), x\rangle
= -\conj{\langle x, u(x)\rangle} = -\conj\lambda\,\norm x^2 ,
\]
ومنه $\lambda = -\conj\lambda$: أي تخيلية محضًا. وبما أن $(\iu
u)^* = -\iu\,u^*$ (\omterm{def:b2:hermitian:adjoint}{فالمرافق} مرافق الخطية في الأعداد)،
يعطي $u^* = u$ أن $(\iu u)^* = -\iu u$: أي إن التطبيق $u \mapsto \iu u$
يرسل الهرميتية إلى ضدّ الهرميتية، بمقلوب $w \mapsto -\iu
w$: أي إنه تقابل. والتشاكل ضدّ الهرميتي $u$ يحقق $u^*u = -u^2 =
uu^*$: فهو ناظمي، ومن ثم \omterm{def:b2:reduction:diag}{قابل للتقطير} \omterm{def:b2:hermitian:adjoint}{وحدويًا} حسب
\cref{exo:b2:hermitian:8}.
\end{solution}

\begin{solution}{exo:b2:hermitian:10}
\begin{enumerate}
  \item يبدّل $S$ أساسًا متعامدًا ممنظمًا: $\norm{Sx} = \norm
        x$، ومن ثم فإن $S$ وحدوي. وبترقيم الإحداثيات بالدليل $j = 0,
        \dots, n-1$ بترديد $n$: $(Sx)_j = x_{j-1}$، ومن ثم من أجل
        $(f_k)_j = \frac{\omega^{jk}}{\sqrt n}$:
        \[
        (Sf_k)_j = \frac{\omega^{(j-1)k}}{\sqrt n}
        = \omega^{-k}\,(f_k)_j :
        \qquad Sf_k = \omega^{-k}f_k .
        \]
        والتعامد والتنظيم: $\langle f_k, f_l\rangle = \frac1n
        \sum_j \omega^{j(l-k)} = \delta_{kl}$ (فالمجموع الهندسي
        لجذر وحدة غير محايد ينعدم).
  \item $Cf_k = \sum_j c_j S^jf_k = \bigl(\sum_j
        c_j\omega^{-jk}\bigr)f_k = \widehat c(k)\,f_k$: فكل
        مصفوفة دائرية قطرية في أساس فورييه المتعامد الممنظم،
        ومن ثم فهي ناظمية، \omterm{def:b2:reduction:eigen}{بطيف} $\{\widehat c(k)\}$. (وتغيير
        الأساس هو تحويل فورييه المتقطع:
        فيصير الالتفاف ضربًا.)
\end{enumerate}
\end{solution}

\begin{solution}{exo:b2:hermitian:11}
($\Leftarrow$) ليكن $P^2 = P = P^*$. وينقسم كل $v$ إلى $v = Pv
+ (v - Pv)$ مع $Pv \in \operatorname{im} P$ و$P(v - Pv) =
0$. والقطعتان متعامدتان: فمن أجل أي $x, y$،
\[
\langle Px, (I - P)y\rangle = \langle x, P(I-P)y\rangle
= \langle x, (P - P^2)y\rangle = 0 :
\]
أي $\ker P \perp \operatorname{im} P$، ومن ثم فإن $P$ هو الإسقاط
المتعامد على صورته. ($\Rightarrow$) وإذا كان $P$ هو
الإسقاط المتعامد على $F = \operatorname{im}P$: فمن أجل كل
$x, y$، $\langle Px, y\rangle = \langle Px, Py\rangle$ (فتسقط
المركّبة $y - Py \perp F$) وبالتناظر $\langle
x, Py\rangle = \langle Px, Py\rangle$: أي $\langle Px, y\rangle =
\langle x, Py\rangle$، أي $P^* = P$. والمصفوفات: على $\C v$
(مع $\norm v = 1$): $Px = v\,\langle v, x\rangle$، أي $P =
vv^\dagger$؛ وعلى $\operatorname{Vect}(v_1, \dots, v_k)$
المتعامد الممنظم: $P = \sum_i v_iv_i^\dagger$.
\end{solution}

\begin{solution}{exo:b2:hermitian:12}
نقطّر $A = U D U^\dagger$ (بالمبرهنة الطيفية)، مع $D$ قطرية
بعناصر من بين $\lambda_i$. عندئذٍ $P_i = L_i(A) = U
L_i(D)U^\dagger$، و$L_i(D)$ قطرية بعناصر
$L_i(\lambda_j) = \delta_{ij}$: أي بواحدات عند خانات
$E_i$ بالضبط. ومن ثم فإن $P_i$ هرميتية (لأن $L_i$ حقيقية و$D$ حقيقي)، ومتساوية القوى،
وصورتها $E_i$ ونواتها $\bigoplus_{j\neq i}E_j = E_i^\perp$
(بتعامد \omterm{def:b2:reduction:eigen}{الفضاءات الذاتية}): أي إنها الإسقاط المتعامد على
$E_i$ (\cref{exo:b2:hermitian:11}). وتعطي الأنماط القطرية المنفصلة
أن $P_iP_j = 0$ (لأن $i \neq j$)؛ و$\sum_i L_i = 1$ (فالدرجة $< p$،
والقيمة $1$ عند $p$ نقطة)، ومنه $\sum P_i = I$؛ ويعطي $\sum_i
\lambda_iL_i(\lambda_j) = \lambda_j$ أن $A = \sum
\lambda_iP_i$. ومن أجل أي كثير حدود $f$: يكون للمصفوفة $f(D)$ القطر
$f(\lambda_j)$، ومنه
\[
f(A) = \sum_{i=1}^{p} f(\lambda_i)\,P_i :
\]
فتُحسب دوال $A$ طيفيًا --- وهو الحساب الذي
يمدّده مجلد السنة الثالثة إلى $f$ المتصلة وما بعدها.
\end{solution}

\begin{solution}{pb:b2:hermitian:1}
\textbf{1.} $\conj{\langle x, Ax\rangle} = \langle Ax, x\rangle
= \langle x, A^*x\rangle = \langle x, Ax\rangle$: فهو حقيقي. وبكتابة
$x = \sum c_ie_i$:
\[
R_A(x) = \frac{\sum_i\lambda_i\abs{c_i}^2}
{\sum_i\abs{c_i}^2} ,
\]
وهو متوسط موزون للقيم الذاتية: فيقع في
$\intcc{\lambda_n}{\lambda_1}$، والحدّان مبلوغان عند
$e_1$ و$e_n$.

\textbf{2.} من أجل $t$ حقيقي وأي $v$، ننشر $R_A(x + tv) =
\frac{N(t)}{D(t)}$ مع
\[
N(t) = \langle x, Ax\rangle + 2t\Re\langle v, Ax\rangle +
t^2\langle v, Av\rangle,
\quad
D(t) = \norm x^2 + 2t\Re\langle v, x\rangle + t^2\norm v^2 .
\]
والمشتق عند $t = 0$ هو
\[
\frac{2}{\norm x^2}\,
\Re\bigl\langle v,\ Ax - R_A(x)\,x\bigr\rangle .
\]
وهو ينعدم من أجل كل $v$ إذا وفقط إذا كان $\Re\langle v, w\rangle = 0$ من أجل
كل $v$، حيث $w = Ax - R_A(x)x$؛ وتعويض $v$ بالمقدار $\iu v$
يقتل الجزء التخيلي أيضًا: أي $w = 0$، أي $Ax = R_A(x)x$.
فالنقاط الحرجة للمقدار $R_A$ هي \omterm{def:b2:reduction:eigen}{المتجهات الذاتية} بالضبط، وقيمتها
الحرجة هي \omterm{def:b2:reduction:eigen}{القيمة الذاتية}.

\textbf{3.} من أجل $x = \sum_{i\leq k}c_ie_i \in V_k$: يكون $R_A(x)$
متوسطًا موزونًا للقيم $\lambda_1, \dots, \lambda_k$، ومن ثم
$\geq \lambda_k$، مع المساواة عند $e_k$: أي $\min_{V_k} R_A =
\lambda_k$. وبالتناظر على $W_k$ يشمل المتوسط
القيم $\lambda_k, \dots, \lambda_n$: أي $\max_{W_k}R_A = \lambda_k$.

\textbf{4.} ليكن $\dim V = k$. عندئذٍ $\dim V + \dim W_k = n + 1 >
n$، ومن ثم توجد متجهة واحدية $x \in V \cap W_k$، ويكون $R_A(x) \leq
\lambda_k$ (السؤال 3): أي $\min_{V}R_A \leq \lambda_k$ من أجل كل
$V$ كهذا. وبما أن $V_k$ يبلغ $\lambda_k$، فإن الأعظمية الصغرى تساوي
$\lambda_k$. وصيغة الأصغرية العظمى هي الحجة نفسها بدورين
معكوسين (فالشرط $\dim W = n - k + 1$ يفرض $W \cap V_k \neq
\{0\}$، ومنه $\max_W R_A \geq \lambda_k$، وهو مبلوغ عند $W_k$).

\textbf{5.} $R_B(x) = R_A(x) + \frac{\langle x,
(B-A)x\rangle}{\norm x^2} \geq R_A(x)$ \omterm{def:b2:funcseq:def}{نقطيًا}. وبأخذ
$\min$ على أي فضاء $V$ بعده $k$ ثم $\max$ على $V$:
$\lambda_k(B) \geq \lambda_k(A)$ حسب السؤال 4.

\textbf{6.} غراسمان: $\dim(V\cap W) = \dim V + \dim W -
\dim(V + W) \geq \dim V + \dim W - n > 0$. وبتطبيق هذا مرتين:
\[
\dim(V_1\cap V_2\cap V_3)
\geq \dim(V_1\cap V_2) + \dim V_3 - n
\geq \dim V_1 + \dim V_2 + \dim V_3 - 2n .
\]

\textbf{7.} للفضاءات الجزئية $W_i(A)$ و$W_j(B)$ (السؤال 3، من أجل
$A$ و$B$) و$V_{i+j-1}(A+B)$ الأبعادُ $(n-i+1) +
(n-j+1) + (i+j-1) = 2n + 1 > 2n$: فيوجد بالسؤال 6 متجهة واحدية
$x$ في الثلاثة معًا. عندئذٍ
\[
\lambda_{i+j-1}(A+B) \leq R_{A+B}(x)
= R_A(x) + R_B(x) \leq \lambda_i(A) + \lambda_j(B),
\]
والمتراجحة اليسرى لأن $x \in V_{i+j-1}(A+B)$ (السؤال
3)، واليمنى بالمقدارين $W$.

\textbf{8.} في أساس طيفي للمصفوفة $E$: $\norm{Ex}^2 = \sum
\lambda_k(E)^2\abs{c_k}^2 \leq \max_k\lambda_k(E)^2\,\norm
x^2$، وهو مبلوغ عند \omterm{def:b2:reduction:eigen}{المتجهة الذاتية} الموافقة: أي $\vertiii E_2
= \max_k\abs{\lambda_k(E)}$. وفايل مع $j = 1$: $\lambda_k(A+E)
\leq \lambda_k(A) + \lambda_1(E) \leq \lambda_k(A) + \vertiii
E_2$؛ وبتطبيق هذا على $(A+E) + (-E)$: $\lambda_k(A) \leq
\lambda_k(A+E) + \vertiii E_2$. ومعًا:
$\abs{\lambda_k(A+E) - \lambda_k(A)} \leq \vertiii E_2$.

\textbf{9.} الحدود الدنيا: $P \geq 0$ والسؤال 5. والعليا: رتبة $P$
هي $1$، ومنه $\lambda_2(P) = 0$؛ وفايل مع $i = k-1$، $j =
2$:
\[
\lambda_k(A + P) \leq \lambda_{k-1}(A) + \lambda_2(P)
= \lambda_{k-1}(A) .
\]

\textbf{10.} $A + E = \begin{pmatrix} 2 & 1+\iu\\ 1-\iu &
2\end{pmatrix}$: \omterm{def:b2:reduction:charpoly}{كثير الحدود المميز} $(2-\lambda)^2 -
\abs{1+\iu}^2 = (2-\lambda)^2 - 2$، \omterm{def:b2:reduction:eigen}{والطيف} $\{2 + \sqrt2,\ 2
- \sqrt2\}$. وفي مقابل $\operatorname{Sp}A = \{3, 1\}$:
\[
\abs{(2+\sqrt2) - 3} = \abs{(2-\sqrt2) - 1} = \sqrt2 - 1
\approx 0.414 \leq 1 = \vertiii E_2 . \checkmark
\]

\textbf{11.} ننظر إلى $\C^{n-1} = \operatorname{Vect}(e_1, \dots,
e_{n-1})$ داخل $\C^n$ (بالأساس القياسي): فمن أجل $x$ هناك،
$\langle x, Bx\rangle = \langle x, Ax\rangle$، ومن ثم فإن $R_B$ هو
مقصور $R_A$. والحدّ الأعلى: تجري الأعظمية الصغرى من أجل $\lambda_k
(B)$ على الفضاءات الجزئية ذوات البعد $k$ \emph{من}
$\C^{n-1}$، وهي عائلة جزئية من فضاءات $\C^n$: أي $\lambda_k(B) \leq
\lambda_k(A)$. والحدّ الأدنى: تجري الأصغرية العظمى من أجل $\lambda_k(B)$
على فضاءات $\C^{n-1}$ الجزئية ذوات البعد $(n-1)-k+1 =
n-k$؛ وكلٌّ منها أيضًا فضاء جزئي من $\C^n$ بعده $n -
(k+1) + 1$، ومن ثم فأعظميته $\geq \lambda_{k+1}(A)$:
أي $\lambda_k(B) \geq \lambda_{k+1}(A)$.

\textbf{12.} نحذف الأسطر والأعمدة واحدًا واحدًا ونسلسل
السؤال 11: فيزيح كل حذف الدليل الأدنى بواحد، فيعطي
$\lambda_{k+m}(A) \leq \lambda_k(B) \leq \lambda_k(A)$.

\textbf{13.} إن مقارنة $A$ بمصفوفة تبديلة (وهي وحدوية)
لا تغيّر \omterm{def:b2:reduction:eigen}{الطيف} ولا مجموعة العناصر القطرية
المتعددة: فلنفترض أن $d_1, \dots, d_k$ تحتل المواضع الرائدة.
وعندئذٍ يكون للمصفوفة الجزئية الرئيسية الرائدة $B$ ذات القياس $k\times k$
الأثرُ $\operatorname{tr} B = \sum_{i\leq k}d_i$، وتحقق قيمها الذاتية
أن $\mu_i(B) \leq \lambda_i(A)$ (السؤال 12): وبالجمع،
$\sum_{i\leq k}d_i \leq \sum_{i\leq k}\lambda_i(A)$. وعند $k =
n$ يساوي الطرفان $\operatorname{tr} A$.

\textbf{14.} يعطي أخذ $x_i = e_i$ القيمةَ
$\sum_{i\leq k}\lambda_i$: فالأعظمية $\geq$. وعكسيًا، تمتدّ
العائلة المتعامدة الممنظمة $(x_1, \dots, x_k)$ إلى
أساس متعامد ممنظم، أي إلى مصفوفة وحدوية $U$ أعمدتها الأولى
$x_i$؛ وعندئذٍ يكون $\sum_i\langle x_i, Ax_i\rangle$ مجموع
العناصر القطرية $k$ الأولى للمصفوفة $U^\dagger AU$، وهو حسب السؤال
13 لا يتجاوز مجموع قيمها الذاتية $k$ الكبرى ---
أي $\sum_{i\leq k}\lambda_i(A)$. ومنه ينتج مبدأ كاي فان
للأعظمية.

\textbf{15.} التداخل ($\operatorname{Sp}A = \{2+\sqrt2, 2,
2-\sqrt2\}$، $\operatorname{Sp}B = \{3, 1\}$):
\[
2 \leq 3 \leq 2 + \sqrt2,
\qquad
2 - \sqrt2 \leq 1 \leq 2 . \checkmark
\]
وشور بقطر $(2,2,2)$: $2 \leq 2+\sqrt2$؛ $4 \leq 4 +
\sqrt2$؛ $6 = 6$ (وهو الأثر). \checkmark

\textbf{16.} $b_{ii} - a_{ii} = \langle e_i, (B-A)e_i\rangle
\geq 0$؛ وبالجمع نحصل على الآثار. ومن أجل أي $C$: $\langle x,
C^\dagger(B - A)Cx\rangle = \langle Cx, (B-A)(Cx)\rangle \geq
0$: أي $C^\dagger AC \leq C^\dagger BC$.

\textbf{17.} $A \geq 0$ واضح؛ و$B - A = \begin{pmatrix} 1 & 1\\
1 & 1\end{pmatrix}$ شبه معرَّفة موجبة (بقيم ذاتية $2,
0$): أي $A \leq B$. لكن
\[
B^2 = \begin{pmatrix} 5 & 3\\ 3 & 2\end{pmatrix},
\qquad
B^2 - A^2 = \begin{pmatrix} 4 & 3\\ 3 & 2\end{pmatrix},
\qquad \det(B^2 - A^2) = -1 < 0 :
\]
ليست شبه معرَّفة موجبة. فالتربيع لا يحترم
ترتيب لوفنر.

\textbf{18.} ليكن $S = \sqrt A$، $T = \sqrt B$ (وهما هرميتيتان
شبه معرَّفتين موجبتين، فالمصفوفة \cref{exo:b2:hermitian:7} ممدَّدة إلى
شبه المعرَّفة بالصيغة الطيفية نفسها). و$T - S$
هرميتية؛ ولتكن $\mu$ أي \omterm{def:b2:reduction:eigen}{قيمة ذاتية} لها، ولتكن $v$ \omterm{def:b2:reduction:eigen}{متجهة ذاتية} واحدية.
ومن $T^2 - S^2 = T(T - S) + (T - S)S$:
\[
0 \leq \langle v, (B - A)v\rangle
= \langle Tv, (T-S)v\rangle + \langle (T-S)v, Sv\rangle
= \mu\bigl(\langle v, Tv\rangle + \langle v,
Sv\rangle\bigr)
\]
(لأن $\mu$ حقيقي). فإذا كان $\langle v, Tv\rangle + \langle v,
Sv\rangle > 0$، فإن $\mu \geq 0$. وإذا انعدم، فإن الحدّين
غير السالبين ينعدمان معًا؛ ويفرض $\langle v, Tv\rangle =
\norm{T^{1/2}v}^2 = 0$ أن $Tv = 0$، وكذلك $Sv = 0$، ومنه
$\mu v = (T - S)v = 0$ و$\mu = 0$. ومن ثم فكل قيم $T -
S$ الذاتية $\geq 0$: أي $\sqrt A \leq \sqrt B$.

\textbf{19.} التوافق بالمصفوفة $A^{-1/2}$ (السؤال 16): $I \leq M
:= A^{-1/2}BA^{-1/2}$. ومن ثم فكل قيم $M$ الذاتية $\geq 1$،
ومنه تقع قيم $M^{-1}$ الذاتية في $\intoc{0}{1}$: أي $M^{-1} \leq I$.
لكن $M^{-1} = A^{1/2}B^{-1}A^{1/2}$؛ ويعطي توافق $M^{-1} \leq
I$ بالمصفوفة $A^{-1/2}$ أن $B^{-1} \leq A^{-1}$.

\textbf{20.} $\sqrt A^{-1}(AB)\sqrt A = \sqrt A\,B\sqrt A$: ومن ثم فإن
$AB$ مشابهة للمصفوفة الهرميتية المعرَّفة الموجبة $\sqrt
A\,B\sqrt A$ (وهي معرَّفة: لأن $\langle x, \sqrt AB\sqrt Ax\rangle =
\langle \sqrt Ax, B\sqrt Ax\rangle > 0$): فقيمها الذاتية
حقيقية وموجبة. زيادةً على ذلك
\[
\langle x, \sqrt AB\sqrt Ax\rangle
\leq \lambda_{\max}(B)\,\norm{\sqrt Ax}^2
= \lambda_{\max}(B)\,\langle x, Ax\rangle
\leq \lambda_{\max}(A)\lambda_{\max}(B)\norm x^2 ,
\]
ومنه $\lambda_{\max}(AB) = \max R_{\sqrt AB\sqrt A} \leq
\lambda_{\max}(A)\lambda_{\max}(B)$.

\textbf{21.} مع $v_k = (\sin j\theta_k)_j$ والمتطابقة
$\sin((j-1)\theta) + \sin((j+1)\theta) =
2\sin(j\theta)\cos\theta$: من أجل $2 \leq j \leq n-1$،
\[
(T_nv_k)_j = -\sin((j{-}1)\theta_k) + 2\sin(j\theta_k) -
\sin((j{+}1)\theta_k)
= (2 - 2\cos\theta_k)\sin(j\theta_k) .
\]
ويعمل السطر $1$ لأن $\sin(0\cdot\theta_k) = 0$، ويعمل السطر $n$
لأن $\sin((n+1)\theta_k) = \sin(k\pi) = 0$: فتختار شروط الحدّ
$\theta_k = \frac{k\pi}{n+1}$ بالضبط. ومن ثم
$T_nv_k = 4\sin^2\bigl(\frac{k\pi}{2(n+1)}\bigr)v_k$؛ والقيم $n$
مختلفة في $\intoo04$ وكذلك $v_k \neq 0$: فهذا هو
\omterm{def:b2:reduction:eigen}{الطيف} كله، وهو موجب، ومن ثم فإن $T_n$ معرَّفة موجبة.

\textbf{22.} $\min_id_i\,I \leq D \leq \max_id_i\,I$، ومنه $T_n +
\min_id_i\,I \leq T_n + D \leq T_n + \max_id_i\,I$
(فإضافة $T_n$ تحفظ الترتيب)؛ ويعطي السؤال 5 مع
$\lambda_k(T_n + cI) = \lambda_k(T_n) + c$ الحصرَ.

\textbf{23.} \omterm{def:b2:reduction:eigen}{طيف} $T_2 = \begin{pmatrix} 2 & -1\\ -1 &
2\end{pmatrix}$ هو $\{3, 1\}$، و
\[
2 - \sqrt2 \;\leq\; 1 \;\leq\; 2 \;\leq\; 3 \;\leq\; 2 +
\sqrt2 :
\]
فتداخل كوشي متحقَّق منه. (وهذه الأطياف هي نفسها التي في
السؤال 15: فالمقارنة بالمصفوفة $\operatorname{diag}(1,-1,1)$ تقلب
إشارة العناصر خارج القطر.) والقراءة البيانية: $T_2$ هي
المصفوفة من نمط لابلاس للمسار بعد حذف الرأس الأخير
--- أي مصفوفة جزئية رئيسية، وهو بالضبط وضع السؤال
11.

\textbf{24.} تسلك القيمتان الذاتيتان الحدّيتان هكذا:
\[
\lambda_{\min}(T_n) = 4\sin^2\frac{\pi}{2(n+1)}
\sim \frac{\pi^2}{(n+1)^2},
\qquad
\lambda_{\max}(T_n) = 4\cos^2\frac{\pi}{2(n+1)}
\longrightarrow 4 .
\]
ومنه
\[
\kappa_n = \frac{\lambda_{\max}}{\lambda_{\min}}
\sim \frac{4(n+1)^2}{\pi^2} :
\]
فكلما دقّت الشبكة ساء شرط المشتق الثاني
المتقطع --- وهي واقعة توجّه تصميم \omterm{def:b2:structures:algebra}{الجبر} الخطي
العددي.

\textbf{25.} (1) قد تتحرك جذور \omterm{def:b2:reduction:charpoly}{كثير الحدود المميز}
تحركًا جامحًا تحت اضطرابات مصفوفة عامة، لكن
توصيف الأصغرية العظمى يثبّت كل \omterm{def:b2:reduction:eigen}{قيمة ذاتية} هرميتية بين
قيم تحسين صريحة، فيفرض استقرار ليبشيتز
بالثابت $1$ في السؤالين 8--9. (2) واستعملت الأسئلة 1 و5 و16--20
قيم رايلي الحدّية وحدها؛ وأما فايل والتداخل وشور
وكاي فان (الأسئلة 7--14) فاحتاجت فعلًا الأصغرية العظمى
الكاملة على الفضاءات الجزئية. (3) ويحوّل ترتيب لوفنر هذه
المتراجحات العددية إلى حساب لمتراجحات المصفوفات --- بفخاخ
حقيقية (السؤال 17) ومبرهنات حقيقية (السؤالان 18--19).
(4) وهذه الأدوات هي الخبز اليومي للتحليل العددي
(مصفوفات الصلابة في السؤال 24)، ولنظرية الاضطراب
الكمومية (فايل: فمستويات الطاقة تتحرك بمقدار لا يتجاوز \omterm{def:b2:nvs:norm}{معيار}
الاضطراب)، ولمبدأ الأصغرية العظمى في مجلد السنة الثالثة
من أجل المؤثرات المتراصة ذاتية \omterm{def:b2:hermitian:adjoint}{المرافق}.
\end{solution}
